\section{结果}
\label{Sec3}

本节应用上面建议的处方，给出 RI$'$ 与 $\MSb$ 两种方案下非微扰 $Z$ 因子的结果。
出于演示目的，我们聚焦于 5 步超立方（HYP）涂抹的数据，这种涂抹压制幂发散，
并使结果更接近重正化的核子矩阵元。我们聚焦于螺旋度准 PDF 的乘性重正化，
只简要讨论非极化算符的情形。

\subsection{RI$'$ 方案及其到 $\MSb$ 方案的转换}

作为出发点，我们对图~\ref{fig:C} 中用过的两个 RI$'$ 标度 $\bar\mu_0$ 的值（``平行''与``对角''）做了计算。
把两种情形都转换到 $\bar\mu{=}2$ GeV 的 $\MSb$ 方案之后，便可以量化与格点伪影以及转换因子截断
相关的系统不确定性，如下一小节所解释。

在图~\ref{fig:ZA} 中，我们给出提取的螺旋度 $Z$ 因子 $Z_{\Delta h}$ 的数值，它重正化裸矩阵元
$\Delta h(P_3,z)$。每幅图中我们都叠加了 RI$'$（空心符号）与 $\MSb$（实心符号）方案的结果，
包括 $Z$ 因子的实部和虚部。我们采用动量源技术~\cite{Gockeler:1998ye,Alexandrou:2015sea}，
它以少量测量（典型为 ${\cal O}(10)$ 次）提供很高的统计精度。如图所示，统计误差几乎不可见。
左（右）图对应 $\bar\mu_0$ 的``平行''（``对角''）选择。%
我们发现 $Z_{\Delta h}^{\overline{\rm MS}}$ 的虚部相比 $Z_{\Delta h}^{\rm RI'}$ 中的对应部分有所减小，
并且在较小的 $z$ 值处相当小。这在图~\ref{fig:ZA} 右图中对``对角''$\bar\mu_0$ 尤为明显：
$\Im[Z_{\Delta h}^{\MSb}]$（蓝色圆圈）小于来自``平行''标度的相应数据，特别是在较大的 $z$ 值处。
值得一提的是，维数正规化和 $\MSb$ 方案中的微扰 $Z$ 因子在微扰论的所有阶都是实的，
因为它仅从极点中提取。因此可以预期，非微扰估计的虚部应当受到强烈压制。
``对角''标度的行为令人鼓舞：直到 ${\sim}10\,a$ 的 $|z|$，其虚部都非常接近于零。
\vspace*{0.25cm}

\begin{figure}[h]
\centering
\includegraphics[scale=0.275,angle=-90]{ZA_RI_MS_Landau_pL3_HYP5_mom_parallel.eps}\,\,\,
\includegraphics[scale=0.275,angle=-90]{ZA_RI_MS_Landau_pt_7_pL3_3_3_HYP5_mom3_diagonal.eps}
\vspace*{-0.5cm}
\begin{minipage}{15cm}
\hspace*{3cm}
\caption{\small{矩阵元 $\Delta h(P_3,z)$（$a P_3=\frac{6\pi}{L}$）的 $z$ 依赖重正化函数。
$\vec{\bar{\mu}}_0$ 的``平行''和``对角''选择分别显示在左图和右图中。
空心（实心）符号对应 RI$'$（$\MSb$）的估计值。}}
\label{fig:ZA}
\end{minipage}
\end{figure}

\vspace*{0.35cm}
如同以前把 RI$'$ 方案应用于提取超局域算符的重正化函数时那样（参见例如文献~\cite{Alexandrou:2015sea}），
这里我们也使用一系列 RI$'$ 重正化标度 $\bar\mu_0$ 的取值。这使我们能够研究顶点函数随外部动量的标度依赖，
并确定能尽可能可靠地提取重正化因子的窗口。与局域流不同的是，${\cal O}(g^2\,a^\infty)$ 格点伪影的计算
极其繁难，目前尚不可得。因此必须小心地选择 $\bar\mu_0$：既要避免非微扰污染，又要满足
$\hat{P}$（式 (\ref{Phat})）较小——最好低于 0.3——的判据，以避免增强的截止效应。
我们使用表~\ref{tab1} 报告的数值计算重正化函数，并从中选出用于拟合的最优范围（采用对角选择）。
平行的 $\bar\mu_0$ 只被用于探索小节~\ref{sub3.2} 中的系统不确定性。

{\small{
\begin{table}[!h]
\vspace*{-0.25cm}
\begin{center}
\begin{tabular}{ccccc}
\hline
\hline\\[-2ex]
标记 & $(n_t,n_x,n_y,n_z)$   &  $(a \bar\mu_0)^2$ &  $\bar\mu$ (GeV)  &  $\hat{P}$  \\[0.5ex]
\hline
\hline\\[-2ex]
& & 对角 & & \\
\hline\\[-2ex]
 m1 &(4,3,3,3)		&1.236       &2.671	 &0.261  		\\
 m2 &(5,3,3,3)		&1.332       &2.773	 &0.251 		\\
 m3 &(6,3,3,3)		&1.448       &2.891	 &0.251		\\
 m4 &(7,3,3,3)		&1.583       &3.023	 &0.261 		\\
 m5 &(8,3,3,3)		&1.737       &3.167	 &0.280		\\
 m6 &(9,3,3,3)		&1.911        &3.321	 &0.306 		\\
 m7 &(10,3,3,3)	&2.104       &3.484 	 &0.339		\\
 m8 &(11,3,3,3)	&2.316       &3.656 	 &0.370 		\\[0.5ex]
 \hline\\[-2ex]
 & & 平行 & & \\
 \hline\\[-2ex]
 m9 &(4,0,0,3)	&0.542       &1.769 	 &0.539 		\\
 m10 &(9,0,0,3)	&1.216       &2.649 	 &0.592 		\\
 m11 &(11,0,0,3)	&1.622       &3.097 	 &0.664 		\\[0.5ex]
\hline
\hline
\end{tabular}
\vspace*{-0.4cm}
\begin{center}
\begin{minipage}{14cm}
\hspace*{3cm}
\caption{\small{RI$'$ 标度的取值，定义为
$a \bar\mu_0{=}\frac{2\pi}{L} (\frac{n_t}{2} {+} \frac{\pi}{4}, n_x,n_y,n_z)$，
其中 $L$ 是格点的空间延展。数值以格点单位和物理单位给出。
最后一列对应式 (\ref{Phat}) 定义的比值 $\hat{P}$。}}
\label{tab1}
\end{minipage}
\end{center}
\end{center}
\vspace*{-0.5cm}
\end{table}
}}

既然我们给出的是螺旋度 PDF 的重正化矩阵元，这一分析聚焦于 $Z_{\Delta h}$。图~\ref{fig:Za_vs_scale} 显示
选定 $z$ 值的 $Z^{\MSb}_{\Delta h}$ 作为 RI$'$ 标度 $(a\,\bar\mu_0)^2$ 的函数，使用标记为 m1$-$m8 的对角值。
左（右）面板为实（虚）部。我们发现对 $(a\,\bar\mu_0)^2$ 存在残余依赖，主要影响虚部。
这是因为顶点函数不仅依赖重正化标度的大小，还依赖其方向——尤其是 $z$ 方向，因为它平行于 Wilson 线。
如果微扰论中的演化已知到更高圈、且离散化效应足够小，那么演化到同一标度后，$Z$ 因子就不应再依赖初始标度。
图~\ref{fig:Za_vs_scale} 所示曲线的非零斜率表明存在不可忽略的格点伪影，同时也是转换因子截断到单圈水平的后果。
然而，要把这两种效应分离开来还需要进一步的研究。我们在小节~\ref{sub3.2} 处理这一问题。
\begin{figure}[!h]
\centering
\includegraphics[scale=0.285,angle=-90]{ZA_B55_MS_HYP5_Re.eps}\,\,\,
\includegraphics[scale=0.285,angle=-90]{ZA_B55_MS_HYP5_Im.eps}
\vspace*{-0.3cm}
\begin{minipage}{15cm}
\hspace*{3cm}
\caption{\small{左：标记为 m1$-$m8 的标度、$z/a{=}5,10,15$ 下 $Z^\MSb_{\Delta h}$ 的实部。
右：同左图，为 $Z^\MSb_{\Delta h}$ 的虚部。}}
\label{fig:Za_vs_scale}
\end{minipage}
\end{figure}

\vspace*{0.15cm}
\noindent
为了消除对 $(a\,\bar\mu_0)^2$ 的残余依赖，我们用函数
\be
\label{eq:extrapol}
Z^{\MSb}_{\Delta h}  = Z_{0, \,{\Delta h}}^{\MSb}  + Z_{1, \,{\Delta h}}^{\MSb} \, .\,(a\,\bar\mu_0)^2
\ee
拟合 $Z^{\MSb}_{\Delta h}$，并把 $Z_{0, \,{\Delta h}}^{\MSb}$ 作为最终结果。该过程对每个 Wilson 线长度 $z$
进行，并对表~\ref{tab1} 给出的 m1$-$m8 对角选择的多个拟合范围重复。对于 2 GeV 处的外推值
$Z_{0, \,{\Delta h}}^{\MSb}$，我们选取在范围 $(a\,\bar\mu_0)^2\,\epsilon\, [1.4,\,2.0]$ 内得到的那个。
排除接近 1 的 $(a\,\bar\mu_0)^2$ 可以避免非微扰效应的污染。上述拟合范围的选择也排除了两个
$\hat{P}{>}0.3$ 的更高标度，它们可能带有可观的格点伪影。作为拟合质量的附加检查，
我们发现所选范围的 $\chi^2/d.o.f$ 很小。图~\ref{fig:ZA_HYP5_extrapolated} 绘出外推值
$Z_{0, \,{\Delta h}}^{\MSb}$（5 步 HYP 涂抹），它将在小节~\ref{sub3.2} 中应用于螺旋度准 PDF 的裸核子矩阵元。

\vspace*{0.5cm}
 \begin{figure}[!h]
\centering
\includegraphics[scale=0.275,angle=-90]{ZA_RI_MS_Landau_extrapolated_HYP5_mom3_diagonal.eps}
\vspace*{-0.3cm}
\begin{minipage}{15cm}
\hspace*{3cm}
\caption{\small{螺旋度算符的外推 $Z$ 因子 $Z_{0, \,{\Delta h}}^{\MSb}$，采用 5 步 HYP 涂抹。
所用的 RI$'$ 标度拟合范围为 $(a\,\bar\mu_0)^2\,\epsilon\, [1.4 {-} 2.0]$。}}
\label{fig:ZA_HYP5_extrapolated}
\end{minipage}
\end{figure}

\vspace*{0.5cm}
获得 $Z_{VV}$ 和 $Z_{VS}$ 的处方也被应用于同一系综。我们发现，HYP 涂抹和``对角''标度的选择
都压制了混合系数 $Z_{VS}$ 和 $Z_{SV}$。图~\ref{fig:Zvv_Zvs} 的左面板绘出标度 ``(7,3,3,3)''
和 5 步 HYP 涂抹下的 $Z^{\rm RI'}_{VV}$ 与 $Z^{\rm RI'}_{VS}$。
$Z^{\rm RI'}_{VS}$ 的实部和虚部大小相当，但至少比乘性因子 $Z^{\rm RI'}_{VV}$ 小一个数量级。
这可以从图~\ref{fig:Zvv_Zvs} 的右面板看出，那里绘出了比值
$\Re[Z^{\rm RI'}_{VS}]/\Re[Z^{\rm RI'}_{VV}]$（红色方块）和
$\Im[Z^{\rm RI'}_{VS}]/\Im[Z^{\rm RI'}_{VV}]$（蓝色圆圈）。从左图可以观察到，
$Z^{\rm RI'}_{VV}$ 和 $Z^{\rm RI'}_{VS}$ 两者的虚部在 $|z/a|<4$ 内都与零一致。
因此，右图中 $|z/a|<4$ 区域内蓝点的大数值并不是显著混合的迹象。
\begin{figure}[!h]
\centering
\includegraphics[scale=0.285,angle=-90]{Zvv_Zvs_RI_Landau_pt_7_pL3_3_3_HYP5_mom3_diagonal.eps}\,\,\,
\includegraphics[scale=0.285,angle=-90]{Zvs_over_Zvv_RI_Landau_pt_7_pL3_3_3_HYP5_mom3_diagonal.eps}
\vspace*{-0.3cm}
\begin{minipage}{15cm}
\hspace*{3cm}
\caption{\small{左：进入非极化准 PDF 重正化的乘性系数（$Z^{\rm RI'}_{VV}$）与混合系数
（$Z^{\rm RI'}_{VS}$），RI$'$ 方案、``(7,3,3,3)'' 标度。右：$Z^{\rm RI'}_{VS}$ 除以
$Z^{\rm RI'}_{VV}$ 的实部之比（红色方块）与虚部之比（蓝色圆圈）。}}
\label{fig:Zvv_Zvs}
\end{minipage}
\end{figure}

\subsection{系统性不确定性的评估}
\label{sub3.2}
\medskip

在本研究的框架中，我们对与单圈转换因子的截断以及离散化效应相关的系统不确定性进行了多项考察。
本小节给出这项研究的主要结论，并给出这些效应的估计。

\vspace*{0.3cm}
在转换到 $\MSb$ 方案之前，我们先对不同 RI$'$ 标度 $\bar\mu_0$ 的取值计算 $Z$ 因子。
例如，图~\ref{fig:ZA} 的两幅图中在标度 m4 和 m9 上给出的 $Z_{\Delta h}^{\rm RI'}$ 是不同的。
如果演化已经收敛，那么演化到同一标度后提取的数值应当一致；这通常发生在微扰论的两圈或三圈。
因此，人们应当能够比较图~\ref{fig:ZA} 中两个标度在 2 GeV 处提取的 $Z_{\Delta h}^\MSb$ 数据。
二者之间的差 $\Delta Z({\rm m4, m9})$ 是存在格点伪影（主要在``平行''情形）、并叠加转换因子截断的一个标志。
图~\ref{fig:ZAdiff} 对实部和虚部展示了这一点；有趣的是，可以看到该差值随 $z$ 变大而增大。
基于这一观察，我们预期``对角''情形也存在格点伪影的这种增大，只是程度较轻。
存在不可忽略系统效应的另一个证据是 $Z_{\Delta h}^\MSb$ 的虚部非零。在图~\ref{fig:ZA} 中，
我们看到``对角''标度有一个小的虚部，并在大 $z$ 值处呈增长趋势。

\vspace*{0.3cm}
\begin{figure}[h!]
\centering
\includegraphics[scale=0.3,angle=-90]{Diff_ZMS_parallel_diagonal.eps}
\vspace*{-0.3cm}
\begin{minipage}{15cm}
\hspace*{3cm}
\caption{\small{$\bar\mu_0$ 的``平行''（m4）与``对角''（m9）情形之间 $Z_{\Delta h}^\MSb$ 的差。
实部和虚部分别用红色方块和蓝色圆圈表示。}}
\label{fig:ZAdiff}
\end{minipage}
\end{figure}

\vspace*{0.5cm}
基于以上论述，我们的结论是：必须处理与格点伪影及转换因子相关的系统不确定性，
才能提取可靠的重正化函数估计。我们打算在不远的将来减小这两种效应，这将消除传播到准 PDF
估计值的系统不确定性。理解支配大 $z$ 区域的不确定性至关重要，因为这些值的矩阵元也进入
产生准 PDF 的傅里叶变换。初步探索表明，两圈贡献的可能量级或许能够压制
$Z_{\Delta h}^{\overline{\rm MS}}$ 的虚部。尽管目前我们还无法精确量化这些系统效应，
但我们将估计其上限。特别地，我们尝试逐个估计每种效应的影响。


%%%%%%%%%%%%%%%%%%%%%%%%%%%%%%%%%%%%%%%%%%%%%
\vskip 0.75cm
\centerline{\bf\underline{\textit{格点伪影}}}
\vskip 0.25cm
%%%%%%%%%%%%%%%%%%%%%%%%%%%%%%%%%%%%%%%%%%%%%

$z{=}0$ 处螺旋度准 PDF 的 $Z$ 因子退化为局域轴向流的重正化函数 $Z_A$。$Z_A$ 与方案和标度无关，
因此 $z{=}0$ 情形使我们得以研究格点伪影。$Z_A$ 已在文献~\cite{Alexandrou:2015sea} 中对该系综做过评估。
那一计算采用了若干个对角标度，连同一种去除格点伪影的技术，提取值为 $Z_A{=}0.7556(5)$。
在本计算的 $z{=}0$ 处，我们用标度 m9 得到 $Z_A{=}0.8620(15)$，用标度 m4 得到 $Z_A{=}0.7727(2)$。
这是又一个迹象：``平行''重正化标度的格点伪影很大，而``对角''情形较轻、但仍不可忽略。

对任意 $z$ 值都可以进行的另一项检验，是比较分量不同、但 $(a\,\bar\mu_0)^2$ 相同的各标度之间的
$Z_{\Delta h}^{\rm \MSb}$。为此，我们利用标记为 m10 和 m11 的数值，它们分别可与 m1 和 m4 比较。
由于这两对的 $(a\,\bar\mu_0)^2$ 近似相同，可以在 RI$'$ 方案中不做任何演化而直接比较它们
\footnote{我们数值确认了从 m1 到 m10 以及从 m4 到 m11 的转换因子偏离 1 小于
1$\text{\textperthousand}$，故将其忽略。}，其差值便可解释为格点伪影。
我们观察到两对的行为相似；作为演示，图~\ref{fig:ZA_diff_m4_m11} 对每个 $z$ 值绘出如下比值：
\be
\label{eq:D}
D_{\rm Re\,(Im)}(\bar\mu_0,\bar\mu_0')\equiv \frac{Z_{\rm Re\,(Im)}^{\rm RI'}(\bar\mu_0) - Z_{\rm Re\,(Im)}^{\rm RI'}(\bar\mu_0')}{Z_{\rm Re\,(Im)}^{\rm RI'}(\bar\mu_0)}\,.
\ee
\begin{figure}[h]
\centering
\includegraphics[scale=0.285,angle=-90]{Diff_ZA_B55_ratio_7333vs11300_RI.eps}
\includegraphics[scale=0.285,angle=-90]{Diff_ZA_B55_ratio_5333vs6333_RI.eps}
\vspace*{-0.3cm}
\begin{minipage}{15cm}
\hspace*{3cm}
\caption{\small{左：使用 $(a\,\bar\mu_0)^2{=}1.583$ 与 $(a\,\bar\mu_0')^2{=}1.622$ 时式 (\ref{eq:D}) 的比值。
右：类似左图，针对两对 $( (a\,\bar\mu_0)^2, (a\,\bar\mu_0')^2 ){=}(1.332, 1.448),\, (1.737,1.911)$。}}
\label{fig:ZA_diff_m4_m11}
\end{minipage}
\end{figure}

\vspace*{0.25cm}
左面板显示 $\bar\mu_0$ 与 $\bar\mu_0'$ 分别对应 m4 和 m11 的情形。偏离零可以归因于格点伪影。
我们发现，从 m11 提取的 $Z$ 因子相对从 m4 提取的值，实部偏离 5-10\%、虚部最高达 30\%。
在 $z/a{\sim}8$ 之后，该偏离似乎稳定到实（虚）部的 5\%（10\%）。

注意，图~\ref{fig:ZA_diff_m4_m11} 的左图给出了``平行''情形相对``对角''情形离散化效应的估计。
计算两个相邻``对角''动量之间的 $D_{\rm Re\,(Im)}$ 以理解伪影的变化也会很有意思。
我们选择了两对 (m2,m3) 和 (m5,m6)，发现 $D_{\rm Re}(\rm{m2,m3})$ 和 $D_{\rm Re}(\rm{m5,m6})$
都小于 1\%，而 $D_{\rm Im}(\rm{m2,m3})$ 和 $D_{\rm Im}(\rm{m5,m6})$ 为 6\% 的量级。
这证实了在式 (\ref{eq:extrapol}) 的拟合中排除``平行''动量能显著减小离散化效应。

%%%%%%%%%%%%%%%%%%%%%%%%%%%%%%%%%%%%%%%%%%%%%
%\vskip 0.75cm
\newpage

\centerline{\bf\underline{\textit{截断效应}}}
\vskip 0.25cm
%%%%%%%%%%%%%%%%%%%%%%%%%%%%%%%%%%%%%%%%%%%%%

为了定量评估只知道单圈微扰论的转换因子的影响，我们对螺旋度 $Z$ 因子的实部和虚部构造比值：
\be
\label{eq:R}
R_{\rm Re\,(Im)}^{{\rm RI'}\,(\MSb)}(z, \bar\mu_0,\bar\mu_0';\bar\mu)\equiv
\frac{Z_{\rm Re\,(Im)}^{{\rm RI'}\,(\MSb)}(z,\bar\mu_0;\bar\mu)}{Z_{\rm Re\,(Im)}^{{\rm RI'}\,(\MSb)}(z,\bar\mu_0';\bar\mu)}
\ee
图~\ref{fig:RIvsMS} 中我们绘出 RI$'$ 与 $\MSb$ 情形的式 (\ref{eq:R})。它们在不同
$(a\bar\mu_0)^2$ 值（m1$-$m8）处提取，并利用文献~\cite{Constantinou:2017sej} 的结果微扰演化到
同一个 $2$ GeV 标度。比值总是以最小的``对角''标度 m1 为参照。
$R_{\rm Re\,(Im)}^{\rm RI'}$ 依赖于到 2 GeV 单圈演化中的截断效应，而
$R_{\rm Re\,(Im)}^\MSb$ 受方案转换截断效应的影响。如果没有来自格点伪影和截断效应的污染，
式 (\ref{eq:R}) 在两种方案中、对所有 $(a\bar\mu_0)^2$ 和 $z/a$ 值都应等于 1。
这一认识使我们可以研究截断效应。

\bigskip
\begin{figure}[h]
\centering
\includegraphics[scale=0.285,angle=-90]{ZA_B55_RI2GEV_MS_HYP5_Re.eps}\,\,\,
\includegraphics[scale=0.285,angle=-90]{ZA_B55_RI2GEV_MS_HYP5_Im.eps}
\vspace*{-0.3cm}
\begin{minipage}{15cm}
\hspace*{3cm}
\caption{\small{$z/a{=}5,10,15$ 时式 (\ref{eq:R}) 的比值作为初始 RI$'$ 标度的函数。
所有结果都已演化到 2 GeV，并以 $(a\bar\mu_0')^2{=}1.236$（m1）处的 $Z$ 因子归一。
左：实部。右：虚部。}}
\label{fig:RIvsMS}
\end{minipage}
\end{figure}

可以观察到，虚部对 x 轴所示的初始 RI$'$ 标度的变化更为敏感。为演示这一点，
我们对图~\ref{fig:RIvsMS} 的所有图保持了相同的 y 轴范围。无论 Wilson 线长度 $z$ 如何，
从不同 RI$'$ 标度 $(\bar\mu_0)^2$ 提取的实部比值之间的差别始终很小：在 RI$'$ 方案中至多约 5\%，
在 $\MSb$ 方案中至多 3\%。如上所述，对于 $z{=}0$ 我们可以与文献~\cite{Alexandrou:2015sea}
中高度可靠的数值比较；我们从标度 m1 得到的当前数值高约 2\%。因此我们可以估计，
$Z$ 因子实部的离散化效应典型为 2-5\% 的量级。

$\MSb$ 比值中的差异是格点伪影与转换截断的组合。我们观察到，转换到 $\MSb$ 方案把
$\MSb$ 方案中的差异降至至多 3\%。因此，实部仅由截断造成的影响似乎不超过 2\%。鉴于我们知道
对小 $z$ 值转换截断效应非常小（特别是它在 $z{=}0$ 处为零，那里 $Z_A$ 与方案无关），
我们可以取 0-2\% 作为估计。最后，考虑到截断效应似乎与格点伪影的影响方向相反，
我们估计 $\Re[Z_{\Delta h}^\MSb]$ 中出现的总效应不应超过 5\%。

图~\ref{fig:RIvsMS} 右面板所示虚部比值的情况有些不同。我们的结论是，在 RI$'$ 方案中对中等和大
$z$（$|z|/a\geq5$\,\footnote{对更小的 $z$ 值，虚部的绝对值很小，式 (\ref{eq:R}) 变得没有意义，
即 $Z_{\Delta h}^\MSb$ 虚部即使很小的绝对变化也可能意味着 $R$ 的巨大变化。}），
来自格点伪影的不确定性可达 10\% 的量级。由图~\ref{fig:RIvsMS} 的右面板，$R_{\rm Im}^\MSb$
的总不确定度对大的 Wilson 线长度可增强到 40\% 左右，对中间长度最高可达 80\%。
此外，$Z_{\Delta h}^\MSb$ 虚部的总不确定度基本上与其自身的量级同阶，因为它预期应为实的。
注意，忽略虚部并不能提供任何解决方案，因为要声称计算完全可靠，虚部就必须从计算中自然得出为零。
于是，转换截断效应成为 $Z$ 因子虚部最主要的不确定性来源。

\bigskip
基于本小节给出的研究，我们在表~\ref{tab2} 中总结了对 $Z_{\Delta h}^\MSb$ 中系统不确定性的定量估计。
总效应的估计不是各孤立效应的简单相加，而是考虑了上文讨论过的它们的符号差异。
对虚部而言，估计针对 $|z/a|\geq5$，因为更小的 $z$ 值时 $\Im[Z_{\Delta h}^\MSb]{\sim}0$，
相对效应便失去意义。

我们想强调，由于 $Z$ 因子与裸矩阵元之间是复数相乘，$Z$ 因子虚部的大不确定度意味着
重正化矩阵元的实部也受到影响。此外，我们所考虑的 $Z$ 因子不确定性对不同 Wilson 线长度
向重正化矩阵元不确定性的传递方式不同，因此我们把这一影响的讨论推迟到下一小节。

{\small{
\begin{table}[!h]
\vspace*{0.25cm}
\begin{center}
\begin{tabular}{ccc}
\hline
\hline\\[-2ex]
效应 & $\Re[Z_{\Delta h}^\MSb]$   &  $\Im[Z_{\Delta h}^\MSb]$ \\[0.5ex]
\hline
\hline\\[-2ex]
格点伪影 & 2-5\% & $\lesssim10$\% \\[1ex]
演化截断& 1-2\% & 1-2\% \\[1ex]
转换截断& $\lesssim2$\% & $\lesssim100$\% \\[2ex]
\hline\\[-1.5ex]
总计 & 3-5\% &  $\lesssim100$\% \\[1ex]
\hline
\hline
\end{tabular}
\vspace*{-0.4cm}
\begin{center}
\begin{minipage}{14cm}
\hspace*{3cm}
\caption{\small{$Z_{\Delta h}^\MSb$ 实部和虚部系统不确定性的定量估计。}}
\label{tab2}
\end{minipage}
\end{center}
\end{center}
\end{table}
}}

\subsection{重正化的结果}
\label{sub3.2}

一旦得到非极化、螺旋度和横向性准 PDF 的 $Z$ 因子，就可以着手把重正化应用到核子矩阵元。
这里我们主要关注螺旋度情形，因为其重正化是乘性的。非极化准 PDF 的情形在本小节末尾简要提及。

图~\ref{fig:deltah} 显示重正化的螺旋度核子矩阵元，并把它们与裸值比较。这是一项直截了当的手续，
因为轴算符没有混合，所以重正化是纯乘性的：
\be
\Delta h^{\MSb} (z) = Z_{\Delta h}^\MSb(z)\cdot \Delta h^{bare}(z)\,\,.
\ee
上面的公式涉及复数量，因此 $\Delta h^{\MSb} (z)$ 是裸矩阵元实部和虚部的混合。不过，
每个 $z$ 值都是独立重正化的。我们使用式 (\ref{eq:extrapol}) 在范围
$(a\,\bar\mu_0)^2\,\epsilon\, [1.4,\,2.0]$ 内外推得到的 $Z_{\Delta h}^\MSb$，如图~\ref{fig:ZA_HYP5_extrapolated} 所示。
从图~\ref{fig:deltah} 可以看到，在小 $z$ 值处，重正化实部相对裸值略有压制
（$\Re [Z_{\Delta h}^\MSb<1]$）。$\Re[\Delta h^{\MSb}]$ 在 $|z/a|{>}8$ 时与零一致，
但统计误差增大。在大 $z$ 区域，重正化对矩阵元虚部的影响是深刻的：与裸值相比，
我们观察到其数值的放大以及最大值向更大 z 的移动。
\bigskip
\begin{figure}[h]
\centering
\includegraphics[scale=0.285,angle=-90]{deltah_renorm_pL3_HYP5_Re.eps}\,\,\,
\includegraphics[scale=0.285,angle=-90]{deltah_renorm_pL3_HYP5_Im.eps}
\vspace*{-0.3cm}
\begin{minipage}{15cm}
\hspace*{3cm}
\caption{\small{$\bar\mu{=}$2 GeV 的 $\MSb$ 方案下螺旋度准 PDF 的重正化矩阵元，
使用从拟合范围 $(a\,\bar\mu_0)^2\,\epsilon\, [1.4,\,2.0]$ 提取的 $Z_{0, \Delta h}^{\MSb}$。}}
\label{fig:deltah}
\end{minipage}
\end{figure}

\vspace*{0.5cm}
螺旋度准 PDF 的重正化矩阵元在此首次给出，它们展示了准 PDF 领域的巨大进展。
然而，在尝试与物理 PDF 比较之前，我们必须理解 $\Delta h^{\MSb} (z) $ 从其重正化函数继承的不确定性。
正如上一小节所论证的，一项稳健的计算需要扣除 $\mathcal{O}(g^2\,a^\infty)$ 格点伪影，
并大幅减少 RI$'$ 与 $\MSb$ 方案间转换的截断不确定性。我们尝试给这些系统效应设定界限，
首先看重正化矩阵元的实部，它读作：
\be
\label{eq:Re}
\Re[\Delta h^\MSb]=\Re[Z_{\Delta h}^{\MSb}]\,\, \Re[\Delta h^{bare}]-\Im[Z_{\Delta h}^{\MSb}]\,\, \Im[\Delta h^{bare}]\,.
\ee
对小 $z$ 值，$\Im[Z_{\Delta h}^{\MSb}]$ 近似为零，因此 $\Re[\Delta h^\MSb]$ 近似等于式 (\ref{eq:Re})
的第一项。相反，当 $|z/a|\gtrsim10$ 时，$\Re[\Delta h^\MSb]$ 从虚部获得显著贡献，
因为 $\Re[\Delta h^{bare}]<\Im[\Delta h^{bare}]$。


重正化矩阵元小 $z$ 区域的不确定性由 $Z$ 因子实部的不确定性主导——如上一小节所论证的，
后者不应超过 5\%。在 $z{=}0$ 观察到的局部极小很可能源于格点伪影和截断效应。
在大 $z$ 区域，若 $Z_{\Delta h}^{\MSb}$ 具有非零虚部、且 $\Im[\Delta h^{bare}]$ 比
$\Re[\Delta h^{bare}]$ 更慢地衰减到零——这正是我们在数据中观察到的——则重正化矩阵元的实部会
从裸矩阵元的虚部接收负的贡献。该效应是非物理的，预期通过把我们的计算扩展到转换的两圈阶
并扣除格点伪影而被强烈压制或消除。

重正化矩阵元的虚部
\be
\label{eq:Im}
\Im[\Delta h^\MSb]=\Re[Z_{\Delta h}^{\MSb}]\,\, \Im[\Delta h^{bare}]+\Im[Z_{\Delta h}^{\MSb}]\,\, \Re[\Delta h^{bare}],
\ee
相比裸矩阵元在中间 $z$ 区间得到增强。这主要是由于 $\Re[Z_{\Delta h}^{\MSb}]$ 随 $z$ 增大的速率
快于 $\Im[\Delta h^{bare}]$ 的衰减。不过，所得数值也接收式 (\ref{eq:Im}) 第二项的贡献，
其中 $\Im[Z_{\Delta h}^{\MSb}]$ 如上一小节所讨论承受着很大的不确定度。

我们想强调，无法对 $\Delta h^\MSb$ 实部和虚部的相对不确定度给出单一数字。
根据式 (\ref{eq:Re}) - (\ref{eq:Im})，不同的 $z$ 区域以不同的方式受 $Z$ 因子实部和虚部的影响。
在小 $z$ 处，$\Im[Z_{\Delta h}^{\MSb}]$ 很小，传播到矩阵元的不确定性由 $\Re[Z_{\Delta h}^{\MSb}]$
的不确定度主导，后者为 5\% 的量级，与当前达到的统计不确定度相当。因此 $\Delta h^\MSb$
在该区域相当稳固。然而当 $z$ 增大时，来自 $\Im[Z_{\Delta h}^{\MSb}]$ 的不确定性开始主导，
并在大 $z$ 时达到 100\%。这样一来，微扰扣除 $\mathcal{O}(g^2\,a^\infty)$ 格点伪影以及扩展到
两圈微扰转换到 $\MSb$ 方案所带来的改进，对获得有意义的重正化矩阵元数值、进而获得有意义的
准 PDF 数值都非常重要。

\vspace*{0.5cm}
矢量算符乘性重正化因子和混合系数的数值可用于通过式 (\ref{h_R}) 正确重正化非极化准 PDF。
成功的重正化需要标量和矢量算符的裸核子矩阵元，以便提取重正化的 $h_V^{\rm RI'}$。
然后必须把它乘以转换因子 $C_V$，把结果带到 $\MSb$ 方案。注意，一旦混合得到处理，
转换因子就是乘性的而非 $2{\times}2$ 矩阵。这是因为连续的维数正规化中没有混合。


\subsection{与光前 PDF 的匹配}
有了重正化矩阵元之后，就可以进行傅里叶变换并得到重正化的准 PDF，它表示以有限动量沿选定空间
（$z$）方向运动的核子中夸克动量的分布：
\begin{equation}
\label{qpdf}
\tilde{q}(x,\mu,P_3) = \int_{-\infty}^\infty \frac{dz}{4\pi} e^{-i z x P_3} \langle N |
\bar{\psi}(0,z)\Gamma W(z) \psi(0,0) |N\rangle^{\MSb,\mu}\,,
\end{equation}
其中 $x$ 是夸克动量分数。在我们的情形中以 $\mu{=}2$ GeV 的 $\MSb$ 方案表达的准 PDF，
随后可以用一圈微扰匹配与同一方案、同一标度下的光前 PDF 相联系。
匹配手续利用了这样一个事实：准 PDF（$\tilde{q}$）与光前 PDF（$q$）只有紫外物理不同 \cite{Xiong:2013bka}。
因此，两者之间的一圈差别表现为有限动量情形与无穷动量情形的顶点修正（下面记作 $Z^{(1)}$）
与波函数修正（$\delta Z^{(1)}$）之差。准 PDF 的一般匹配公式为：\cite{Ji:2015jwa,Chen:2016fxx}
\begin{eqnarray}
\label{invq}
q\left(x,\mu\right) &=& \tilde{q}(x,\mu,P_3) - \frac{\alpha_s(\mu)}{2\pi} \,\tilde{q}(x,\mu,P_3)\,\delta Z^{(1)}\!\left( \frac{\mu}{P_3}\right)  \nonumber\\*
& & - \frac{\alpha_s(\mu)}{2\pi} \int_{-\infty}^{\infty} Z^{(1)}\!\left( \xi, \frac{\mu}{P_3} \right) \tilde{q}\! \left( \frac{x}{\xi},\mu,P_3 \right) \frac{d\xi}{|\xi |} + \mathcal{O}(\alpha_s^2),
\end{eqnarray}
其中 $\alpha_s(\mu)$ 是标度 $\mu$ 处的强耦合常数。在该积分中，我们排除了 $\xi=0$ 附近的小区域，
使得 $\tilde{q}$ 的自变量不至于太大，否则我们会拾取来自式 (\ref{qpdf}) 傅里叶变换周期性导致的
$\tilde{q}$ 镜像的贡献。文献~\cite{Xiong:2013bka} 匹配公式中 $\propto\Lambda/P_3$
（$\Lambda$：横动量截断）的线性发散项不出现在我们的重正化结果中。

\bigskip
\begin{figure}[h]
\centering
\includegraphics[scale=0.55,angle=-90]{qmatched_bare_vs_fit69v2.eps}\,\,\,
\vspace*{-0.3cm}
\begin{minipage}{15cm}
\hspace*{3cm}
\caption{\small{匹配螺旋度 PDF 的比较：一者由完全重正化矩阵元计算的准 PDF 得到（蓝色），
另一者由裸矩阵元乘以局域（$z{=}0$）轴流 $Z$ 因子 $Z_A$ 得到（品红）。仅为定向参考之用，
我们也画出了唯象 PDF（DSSV08 \cite{deFlorian:2009vb} 与 JAM15 \cite{Sato:2016tuz}）。
但我们强调，无意与我们的结果做定量比较，因为仍需仔细考虑一系列系统效应，包括：
截止效应、非物理的 $\pi$ 介子质量、有限体积效应、可能的激发态污染、向无穷核子动量的外推，
以及上一小节所主张的 $\MSb$ 重正化函数计算的改进。}}
\label{fig:matched}
\end{minipage}
\end{figure}

\medskip
在图~\ref{fig:matched} 中，我们展示匹配的螺旋度 PDF：一者使用本工作得到的完全重正化矩阵元（蓝色），
另一者使用裸矩阵元乘以局域（$z{=}0$）轴矢量流重正化函数 $Z_A$（品红），后者对应我们文献
~\cite{Alexandrou:2016jqi} 的结果。我们观察到，本工作的重正化矩阵元朝唯象 PDF 移动，这是有希望的。
特别地，反夸克不对称不再被高估，实际上在当前不确定度下这一不对称变得与零相容。
在夸克部分，匹配后的 PDF 对所有 $x$ 值都有所增强。我们再次强调，与唯象 PDF 的比较只应理解为定性的。
要做定量比较，仍需对一系列系统效应进行仔细研究，包括：截止效应、非物理的 $\pi$ 介子质量、
有限体积效应、可能的激发态污染、向无穷核子动量的外推，以及上一小节所主张的 $\MSb$ 重正化函数
计算的改进：扣除由格点微扰论计算的格点伪影，以及减少微扰转换到 $\MSb$ 方案中的截断效应。
